\chapter*{Introduction Générale}
\addcontentsline{toc}{chapter}{Introduction Générale}

Travailler ou étudier plusieurs heures d’affilée sans se rendre compte qu’on a perdu le fil ; c’est une situation que beaucoup connaissent. La fatigue s’installe progressivement, la posture se dégrade, le téléphone attire l’attention, et pourtant l’utilisateur continue sans vraiment mesurer l’impact de tout ça sur la qualité de son travail. Ce constat, apparemment banal, est à l’origine du projet \textbf{Smart Focus}.

Aujourd’hui, les outils de productivité sont nombreux. Les applications Pomodoro, les bloqueurs de sites, les minuteries ; tout ça existe depuis longtemps. Mais ces outils ont un point commun : ils agissent sur le temps, pas sur l’état réel de l’utilisateur. Ils ne savent pas si vous êtes fatigué, si votre dos est mal positionné, ou si votre regard dérive toutes les deux minutes vers l’écran de votre téléphone. Ils font confiance à l’utilisateur pour s’autoréguler, alors que c’est précisément là que le problème se pose.

C’est dans ce contexte que nous avons conçu Smart Focus. L’idée était de construire un système capable de surveiller passivement le comportement de l’utilisateur pendant une session de travail, d’en extraire des indicateurs concrets : posture, niveau d’éveil, attention, distractions, et de lui fournir un retour utile en temps réel. Le tout sans que l’utilisateur ait à faire quoi que ce soit de particulier.

Le système s’appuie sur un boîtier embarqué équipé d’une caméra et de capteurs, couplé à un backend intelligent et à une application mobile. La vision par ordinateur analyse le flux vidéo, l’intelligence artificielle interprète les résultats, et l’IoT assure le lien entre le matériel et les services numériques. En plus du monitoring, Smart Focus propose un assistant pédagogique : chatbot RAG, génération de quiz, flashcards et planning adaptatif, pour aller au-delà du simple suivi.

Pour mener ce projet, nous avons choisi de travailler en mode agile avec Scrum, ce qui nous a permis de livrer des fonctionnalités progressivement et d’adapter notre trajectoire en fonction des résultats obtenus à chaque sprint.

Ce mémoire retrace l’ensemble du travail réalisé, organisé de la façon suivante :

\begin{itemize}
    \item \textbf{Chapitre 1 : Cadre général du projet} \\
    Ce chapitre présente le contexte dans lequel s’inscrit le projet, la problématique identifiée, et une analyse des solutions existantes qui a orienté nos choix de conception.

    \item \textbf{Chapitre 2 : Sprint 1 – Authentification et interfaces de base} \\
    Nous y décrivons la mise en place du module d’authentification : inscription, connexion, gestion du profil, ainsi que les premières interfaces de l’application mobile.

    \item \textbf{Chapitre 3 : Sprint 2 – Surveillance intelligente de la concentration} \\
    Ce chapitre détaille le pipeline de vision par ordinateur : comment le système analyse le flux vidéo, détecte la fatigue et les distractions, évalue la posture, et calcule un score de concentration global.

    \item \textbf{Chapitre 4 : Sprint 3 – Planning intelligent et assistant d’apprentissage} \\
    Nous présentons ici les modules IA du projet : le chatbot RAG pour interroger des documents PDF, le générateur de quiz et de flashcards, et le système de planification adaptative.

    \item \textbf{Chapitre 5 : Sprint 4 – Système IoT et alertes intelligentes} \\
    Ce dernier chapitre couvre l’intégration du boîtier matériel, le suivi en temps réel depuis l’application mobile, et le système d’alertes basé sur les données de concentration et de sommeil.
\end{itemize}

À travers ce document, nous cherchons à rendre compte non seulement de ce qui a été réalisé, mais aussi des choix qui ont été faits, des problèmes rencontrés, et des compromis acceptés pour aboutir à un système fonctionnel.

\clearpage

%%%%%%%%%%%%%%%%%%%%%%%%%%%%%%